\documentclass[11pt]{amsart}
\usepackage[T1]{fontenc}
\usepackage[utf8]{inputenc}
\usepackage{lmodern,amsmath,amssymb,amsthm,mathtools,mathrsfs}
\usepackage[margin=29mm]{geometry}
\usepackage{microtype,booktabs}
\usepackage[unicode,hidelinks]{hyperref}
\setlength{\emergencystretch}{3em}
\allowdisplaybreaks[2]
\newtheorem{theorem}{Theorem}[section]
\newtheorem{lemma}[theorem]{Lemma}
\newtheorem{proposition}[theorem]{Proposition}
\newtheorem{corollary}[theorem]{Corollary}
\theoremstyle{definition}\newtheorem{definition}[theorem]{Definition}
\theoremstyle{remark}\newtheorem{remark}[theorem]{Remark}
\newcommand{\R}{\mathbb R}\newcommand{\E}{\mathbb E}
\newcommand{\Prob}{\mathbb P}\newcommand{\one}{\mathbf1}
\newcommand{\upref}[2]{#1 in the retained manuscript (\nolinkurl{#2})}
\title[Occupation measures and finite-memory prediction]{Occupation measures in finite-memory prediction}
\author{Qian Qi}
\date{September 8, 2026}
\subjclass[2020]{62B05, 62C10, 90C25, 90C35}
\keywords{attainable information, occupation measures, adaptive acquisition, finite-memory prediction, small-ball bounds, convex duality}
\hypersetup{pdftitle={Occupation measures in finite-memory prediction: A1 v29 new proofs},pdfauthor={Qian Qi}}
\begin{document}
\begin{abstract}
Local small-ball bounds for predictions along a finite acquisition path give a convex variational converse when the occupation law is retained at every loss threshold and checkpoint. We characterize the sharp consequence of those local random-path data. The requirement that the first acquisition precede all observations yields a stronger perspective-flow problem. Its dual provides certified lower values through capacitated minimum-cost flows. Both converses quantitatively dominate the visited-set separator integral without interchanging a maximum and an expectation. In a positive two-trial experiment the resulting all-budget coefficient is $1/8306688$; its gap from the sharp high-resolution average coefficient is less than $4.06$.
\end{abstract}
\maketitle

\section{The revision and its mathematical setting}
This document contains the new proof component of A1 v29. The repository entry point \texttt{main.tex} includes the complete retained scalar, collision, causal-memory, adaptive-graph, phase and policy-menu developments from A1 v26--v28, including the localized growing-graph results, and then incorporates the sections below. The companion manuscript remains a separate, unchanged source. This separately compiled component is not the complete two-volume manuscript.

The controlling report is the independent A1 v26 assessment at commit \texttt{19fbf4fe0e74}. The integration base is the native localized v28 revision at \texttt{090f9875d71a}. The occupation-flow draft was recovered from earlier unpublished work and is integrated here without replacing that localized revision. References to retained source labels below identify dependencies in that manuscript, rather than inventing theorem numbers from an uncompiled volume. Appendix~\ref{app:v29-density} records the acquired-law calculation used in the explicit comparison.

In the statistical application, a finite graph indexes independent positive experiments. At the first visited endpoint of an edge one acquires its first observation block; the other endpoint completes it. A vertex is an atomic input batch. Its commands are independent of the past and revealed only after the vertex is chosen. The persistent state at each boundary has at most $M$ labels. The query descriptor identifies the visited set and one member of a finite future-probe menu, but reveals no discarded continuous history. There is one queried checkpoint per physical run and every trial counts.

For one common controller, write $\mathcal L_j$ for the squared-prediction excess at boundary $j$, averaged over the query menu. Independent coding coins are included in $\xi$ and the potential observation tapes in $\tau$. The average criterion is
\[
 \mathcal R^{\mathrm{av}}_{G,M}
   =\inf_{\mathcal F}\max_{1\le j<|V|}
                         \E_{\tau,\xi}\mathcal L_j.
\]
The maximum-tape criterion replaces $\E_\tau$ by $\sup_\tau$ after averaging $\xi$. Neither criterion chooses separate controllers at separate boundaries. Prediction clipping ensures $0\le\mathcal L_j\le1$. The quantitative local bounds hold after fixing independent seeds, before restricting to a selected visited set. This order of operations is the input to the new argument.

\section{Occupation measures and local prediction bounds}
\label{sec:v29-occupation}
A separator records that a path must visit a costly state. It does not
record the probability with which the other states are visited, and it
permits a different separating set at each loss threshold. We retain the
occupation measure of a single path law throughout the integration. This
gives a convex variational bound adapted to the maximum of the expected
checkpoint losses, rather than to the expectation of their maximum.

\subsection{The finite path experiment}
Let $\mathscr D$ be a finite directed acyclic graph with source $o$, sink
$d$, and internal vertex levels $\mathscr V_1,\ldots,\mathscr V_J$.
Every source--sink path visits exactly one vertex of each level, in that
order. Vertices and arcs on no such path are discarded. Put
$\mathscr V=\bigcup_{j=1}^J\mathscr V_j$ and let $j(w)$ denote the level
of $w$. A random path $P$ is selected on a probability space carrying an
event $A$ of probability $\beta>0$. At level $j$ there is a nonnegative
loss $L_j\le1$. For the graph experiment above, $J=v-1$, $w$ is a visited
set, and $L_j$ is the query-averaged squared-prediction excess. Predictions
may be clipped to $[0,1]$ without increasing any of these losses.

For each $w\in\mathscr V$ suppose that a nondecreasing right-continuous
function $c_w:[0,1]\to[0,1]$ satisfies
\begin{equation}\label{eq:v29-local}
 \Pr\{A,\ w\in P,\ L_{j(w)}\le t\}\le\beta c_w(t)
                         \qquad(0\le t\le1).
\end{equation}
No distribution conditional on the selected vertex is prescribed. In
particular, \eqref{eq:v29-local} is not a conditional independence
hypothesis. The node capacities in
\upref{equation}{eq:v26-node-capacity} give precisely these inequalities.

Write $\mathcal F(c)$ for the polytope of internal vertex occupations
of unit nonnegative $o$--$d$ flows, with the additional constraints
\begin{equation}\label{eq:v29-terminal-cap}
                         0\le x_w\le c_w(1).
\end{equation}
Thus $x_w$ is the incoming, equivalently outgoing, flow at $w$;
conservation is imposed at every internal vertex. The extra constraints
use the bounded range of the loss and are part of the variational
problem. Define
\begin{equation}\label{eq:v29-phi}
 \phi_w(x)=\int_0^1[x-c_w(t)]_+\,dt\quad(0\le x\le1),\qquad
 \Gamma(c)=\min_{x\in\mathcal F(c)}
             \max_{1\le j\le J}\sum_{w\in\mathscr V_j}\phi_w(x_w).
\end{equation}
Every $\phi_w$ is nonnegative, nondecreasing, convex, continuous, and
$1$-Lipschitz. If a path experiment satisfying
\eqref{eq:v29-local} exists, $\mathcal F(c)$ is nonempty. Since the
network is finite, it is compact and the minimum in
\eqref{eq:v29-phi} is attained. This is a finite convex optimization
problem, not a minimization over separately chosen checkpoint encoders.

\begin{theorem}[Occupation-measure converse]\label{thm:v29-occupation}
Every path experiment satisfying \eqref{eq:v29-local} obeys
\begin{equation}\label{eq:v29-occupation-risk}
             \max_{1\le j\le J}\E L_j\ge\beta\Gamma(c).
\end{equation}
More precisely, its single conditional occupation measure
$x_w=\Pr(w\in P\mid A)$ belongs to $\mathcal F(c)$ and satisfies,
simultaneously at all levels,
\begin{equation}\label{eq:v29-level-risk}
 \E[\one_A L_j]\ge
                  \beta\sum_{w\in\mathscr V_j}\phi_w(x_w).
\end{equation}
\end{theorem}
\begin{proof}
The conditional expected arc incidences of $P$ form a unit flow.
At $t=1$, the loss bound and \eqref{eq:v29-local} give
$x_w\le c_w(1)$, so its occupations belong to $\mathcal F(c)$.
Put
\[
 G_w(t)=\beta^{-1}\Pr\{A,\ w\in P,\ L_{j(w)}\le t\}.
\]
Both $G_w(t)\le x_w$ and $G_w(t)\le c_w(t)$ hold. The layer-cake
identity, applied to the subprobability of visiting $w$ on $A$, yields
\[
 \begin{split}
 \E[\one_A\one_{\{w\in P\}}L_{j(w)}]
   &=\beta\int_0^1(x_w-G_w(t))\,dt\\
   &\ge\beta\int_0^1[x_w-c_w(t)]_+\,dt.
 \end{split}
\]
At each level exactly one vertex is visited. Summing at that level
proves \eqref{eq:v29-level-risk}; taking the maximum and then the
minimum over the larger set $\mathcal F(c)$ proves
\eqref{eq:v29-occupation-risk}. The same $x$ was used at every
threshold and every level. No maximum has been moved through an
expectation.
\end{proof}

\begin{proposition}[Comparison with the separator integral]
\label{prop:v29-separator-dominance}
Let
\[
 H(t)=\min_{C\in\operatorname{Sep}(\mathscr D)}
                                      \sum_{w\in C}c_w(t).
\]
Then
\begin{equation}\label{eq:v29-separator-dominance}
                 \Gamma(c)\ge\frac1J\int_0^1[1-H(t)]_+\,dt.
\end{equation}
If the capacities are extended monotonically to $[0,\infty)$,
$H(t)\ge1$ for $t\ge1$ and the upper integration limit may be
replaced by infinity.
\end{proposition}
\begin{proof}
A nonnegative flow in a finite acyclic graph is a convex combination
of path flows: start with a positive source arc, follow positive arcs
to the sink by conservation, subtract the smallest arc value on this
path, and repeat. At least one positive arc is removed on each step.
Consequently every $x\in\mathcal F(c)$ satisfies
$\sum_{w\in C}x_w\ge1$ for every separator $C$; the inequality
allows a path to meet $C$ more than once. Hence
\[
 \sum_{w\in\mathscr V}[x_w-c_w(t)]_+
 \ge\left[1-\min_C\sum_{w\in C}c_w(t)\right]_+.
\]
Integrate, divide by $J$, and use the inequality between the maximum
and the average of the $J$ level sums. Taking the minimum over $x$
gives \eqref{eq:v29-separator-dominance}. At $t=1$, feasibility and
\eqref{eq:v29-terminal-cap} imply $\sum_{w\in C}c_w(1)\ge1$ for
every separator. Monotonicity proves the final assertion.
\end{proof}

\subsection{Power capacities and exactness of the relaxation}
For one positive coordinate allocation of total dimension $k$, write
$c(t)=\min\{1,a t^{k/2}\}$, where $a>0$ includes the label count,
density, recovery norm and determinant product. For every
$0\le x\le c(1)$, direct integration gives
\begin{equation}\label{eq:v29-power-cost}
 \phi(x)=\frac{k}{k+2}\,a^{-2/k}x^{1+2/k}.
\end{equation}
Indeed, the integrand is positive exactly for
$0\le t<(x/a)^{2/k}\le1$. Its integral is
$x(x/a)^{2/k}-a(x/a)^{1+2/k}/(1+k/2)$, which is
\eqref{eq:v29-power-cost}. For a minimum of several positive-allocation
capacities, the integrand is the maximum of their positive integrands;
it is not in general permissible to exchange that maximum with the
integral. Definition~\eqref{eq:v29-phi} retains the exact envelope.
All zero determinant terms remain omitted, as in the preceding
collision-uniform converse.

For a chain with one internal vertex at each of $J$ levels and identical
power capacities, feasibility forces every occupation to be one. If
$a\ge1$, the occupation bound is
\[
              \beta\frac{k}{k+2}a^{-2/k},
\]
whereas the separator integral divided by $J$ is exactly $1/J$ times
this number. This family exhibits an arbitrarily large loss caused
solely by the conversion from a maximum loss to a maximum of expected
losses. It is a statement about the path relaxation; it is not a
size-uniform adaptivity theorem for the original graph experiments.

The variational expression also identifies exactly what can follow
from the local inequalities alone.

\begin{theorem}[Sharp consequence of local path data]
\label{thm:v29-relaxation-exact}
Fix $\mathscr D$, $\beta\in(0,1]$ and the functions $c_w$. Suppose
$\mathcal F(c)\ne\varnothing$. Among all random-path loss models with
$\Pr(A)=\beta$, $0\le L_j\le1$, and
\eqref{eq:v29-local}, the smallest possible value of
$\max_j\E L_j$ is exactly $\beta\Gamma(c)$.
\end{theorem}
\begin{proof}
The lower bound is Theorem~\ref{thm:v29-occupation}. For the opposite
direction choose a minimizing occupation $x$ and decompose its flow
into a distribution on paths. Sample that path independently of an
event $A$ of probability $\beta$. At a vertex $w$ with $x_w>0$,
independently sample $U_w$ uniformly on $[0,x_w]$ and set
\[
 T_w=\inf\{t\in[0,1]:c_w(t)\ge U_w\}.
\]
The set is nonempty almost surely by $x_w\le c_w(1)$.
Monotonicity and right continuity imply
\[
 \Pr(T_w\le t)=\min\{1,c_w(t)/x_w\}.
\]
Put $L_j=T_w$ at the visited vertex of level $j$ on $A$, and put
$L_j=0$ outside $A$. Vertices with $x_w=0$ are never visited and
need no auxiliary variable. Then
\[
 \Pr\{A,w\in P,L_{j(w)}\le t\}
                       =\beta\min\{x_w,c_w(t)\},
\]
and the layer-cake computation in the preceding proof is an equality
at every level. Thus $\max_j\E L_j=\beta\Gamma(c)$.
\end{proof}

The class in Theorem~\ref{thm:v29-relaxation-exact} consists of path
laws and losses satisfying specified local distribution bounds. It does
not require those losses to arise from a common finite-label decoder,
a Bayes experiment or a realizable posterior map. Accordingly, this
exactness statement is not an assertion that the statistical risk in
\upref{equation}{eq:v25-risk-definitions} equals $\beta\Gamma(c)$.
It shows that a stronger consequence requires additional information
not contained in \eqref{eq:v29-local}. The acquisition-time information
constraint supplies such information in the next section.

\section{Information available before the first acquisition}
\label{sec:v29-initial}
The first vertex is chosen before any command or report is revealed.
Conditional on an independent coding seed, that vertex is therefore
fixed. Averaging seeds before recording this fact enlarges the path
relaxation: it allows visits to different first vertices to be selected
as though they could already convey information about the tape. We
retain the seedwise constraint without exchanging the checkpoint maximum
and the expectation.

\subsection{Anchored occupations and perspective costs}
Assume that the underlying tape $\tau$ and a seed $\xi$ have product
law, and that $A$ is a tape event of probability $\beta$. It is enough
that \eqref{eq:v29-local} hold for almost every fixed $\xi$, with the
same functions $c_w$ and the same $\beta$. Suppose also that the first
internal vertex is a function $u(\xi)$, not a function of the tape.
These hypotheses allow the scheduler its complete revealed history
after that first vertex. They also allow coding randomness shared
between scheduler and predictor, provided it was independent of the
tape before acquisition.

For a first-level vertex $u$, let $\mathcal F_u(c)$ be the subset of
$\mathcal F(c)$ carried by paths whose first vertex is $u$.
Discard empty such sets and call the remaining set of anchors $\mathcal U$.
For nonnegative $\alpha_u$ summing to one, choose
$x^u\in\mathcal F_u(c)$ whenever $\alpha_u>0$. Define
\begin{equation}\label{eq:v29-seed-gamma}
 \Gamma_{\rm pre}(c)=
 \min_{\alpha,x^u}\max_{1\le j\le J}
       \sum_{u\in\mathcal U}\alpha_u
                      \sum_{w\in\mathscr V_j}\phi_w(x^u_w).
\end{equation}
This expression has a compact convex formulation. Replace $x^u$ by a
flow $y^u=\alpha_u x^u$ of total mass $\alpha_u$, anchored at $u$, and
impose $y^u_w\le\alpha_u c_w(1)$. Replace each summand by the
perspective
\begin{equation}\label{eq:v29-perspective}
 \mathcal P_w(\alpha,y)=
 \begin{cases}
    \alpha\phi_w(y/\alpha),&\alpha>0,\\
    0,&\alpha=0,\ y=0.
 \end{cases}
\end{equation}
On this feasible domain, the perspective is continuous and convex;
continuity at zero follows from $0\le\phi_w\le1$.
The network conservation laws and the constraints on $(\alpha,y)$ are
linear. Thus \eqref{eq:v29-seed-gamma} is a finite convex program and
attains its minimum.

\begin{theorem}[Pre-acquisition refinement]\label{thm:v29-pre}
Under the seedwise hypotheses above,
\begin{equation}\label{eq:v29-pre-risk}
        \max_j\E_{\tau,\xi}L_j
                \ge\beta\Gamma_{\rm pre}(c)
                \ge\beta\Gamma(c).
\end{equation}
Both bounds apply to a single common controller. No stored command,
report, order prefix or history-correlated random variable is supplied
for free by conditioning on the independent seed.
\end{theorem}
\begin{proof}
For a fixed seed let $x^\xi_w=\Pr_\tau(w\in P\mid A,\xi)$.
Theorem~\ref{thm:v29-occupation}, applied before averaging seeds, gives
at each fixed level
\[
 \E_\tau[\one_A L_j\mid\xi]
                 \ge\beta\sum_{w\in\mathscr V_j}\phi_w(x^\xi_w).
\]
Moreover $x^\xi\in\mathcal F_{u(\xi)}(c)$. Set
$\alpha_u=\Pr_\xi(u(\xi)=u)$ and, for positive $\alpha_u$, set
$x^u=\E_\xi[x^\xi\mid u(\xi)=u]$. Convexity of the anchored flow
polytope gives $x^u\in\mathcal F_u(c)$. Conditional Jensen gives,
for each fixed $j$,
\[
 \E_{\tau,\xi}L_j
 \ge\beta\sum_u\alpha_u\sum_{w\in\mathscr V_j}\phi_w(x^u_w).
\]
Take the maximum only now, and then minimize over the larger feasible
set in \eqref{eq:v29-seed-gamma}. This proves the first inequality.
For the second, $x=\sum_u\alpha_u x^u$ is in $\mathcal F(c)$ and
\[
 \sum_u\alpha_u\sum_{w\in\mathscr V_j}\phi_w(x^u_w)
                   \ge\sum_{w\in\mathscr V_j}\phi_w(x_w)
\]
at every level. Take the maximum and the minimum.
\end{proof}

The refinement is exact for the correspondingly refined abstract path
class. To see this, select a minimizing $(\alpha,x^u)$, let the seed
choose $u$ with probabilities $\alpha_u$, and, conditional on it, use
the construction of Theorem~\ref{thm:v29-relaxation-exact} with flow
$x^u$. All the conditional constructions can be generated from one array of
independent uniform tape variables, independent of the seed, with the
anchor selecting which deterministic sampling map to use. Take $A$
independent of those variables and of the seed with mass $\beta$. The first
vertex is now fixed given the seed, every seedwise capacity holds, and
the level risks are exactly the sums in \eqref{eq:v29-seed-gamma}.
This establishes sharpness within that abstract class, not
attainability by a prescribed statistical experiment.

\subsection{Dual certificates}
A feasible primal flow gives an upper bound on a convex minimum. It
must not be reported as a certified risk lower bound. We give a dual
formula whose feasible values have the required direction.
For $0\le r\le1$ put
\begin{equation}\label{eq:v29-conjugate}
 \widehat\phi_w(r)=\max_{0\le z\le1}\{rz-\phi_w(z)\}.
\end{equation}
The bound on the slopes of $\phi_w$ implies
\begin{equation}\label{eq:v29-support-representation}
 \phi_w(z)=\max_{0\le r\le1}\{rz-\widehat\phi_w(r)\},
                                              \quad 0\le z\le1.
\end{equation}
For a vector $q$ of nonnegative vertex costs define
\[
 d_u(q)=\min_{x\in\mathcal F_u(c)}\sum_w q_wx_w.
\]
This is a capacitated minimum-cost unit-flow problem. When all terminal
capacities $c_w(1)$ equal one, it is the minimum cost of an anchored
source--sink path and is computed by the ordinary acyclic shortest-path
recursion. Otherwise the node capacities must be retained.

Let $\Delta_J=\{\lambda_j\ge0:\sum_j\lambda_j=1\}$. For
$0\le q_w\le\lambda_{j(w)}$ define
\[
 B(\lambda,q)=\sum_w\lambda_{j(w)}
           \widehat\phi_w\bigl(q_w/\lambda_{j(w)}\bigr),
\]
where a term with $\lambda_{j(w)}=q_w=0$ is zero.

\begin{theorem}[Convex-flow duality]\label{thm:v29-dual}
The pre-acquisition variational bound has the exact representation
\begin{equation}\label{eq:v29-dual}
 \Gamma_{\rm pre}(c)=
 \max_{\substack{\lambda\in\Delta_J\\
              0\le q^u_w\le\lambda_{j(w)}\ (u\in\mathcal U)}}
       \min_{u\in\mathcal U}\{d_u(q^u)-B(\lambda,q^u)\}.
\end{equation}
Every feasible choice on the right gives a valid lower certificate.
With a single unanchored flow polytope, the same formula, without the
minimum over $u$, represents $\Gamma(c)$.
\end{theorem}
\begin{proof}
The functions in \eqref{eq:v29-support-representation} are affine in
$z$ and concave in $r$. They are continuous on compact convex sets:
\eqref{eq:v29-conjugate} is a maximum over $0\le z\le1$ and is
Lipschitz in $r$. The classical finite-dimensional convex minimax principle
\cite{Rockafellar1964} therefore gives,
for each fixed $\lambda$ and anchor,
\begin{equation}\label{eq:v29-inner-dual}
 \begin{split}
 m_u(\lambda)
  &:=\min_{x\in\mathcal F_u(c)}\sum_w\lambda_{j(w)}\phi_w(x_w)\\
  &=\max_{0\le q_w\le\lambda_{j(w)}}
                          \{d_u(q)-B(\lambda,q)\}.
 \end{split}
\end{equation}
One may also obtain the inequality from right to left directly by
summing the supporting inequalities
$\lambda_j\phi_w(x_w)\ge q_wx_w-
\lambda_j\widehat\phi_w(q_w/\lambda_j)$ and minimizing over flows.
It is this direction that certifies numerical or exact trial values.

For completeness, the outer minimax step can be seen directly by
separation. Let $K$ be the set of vectors $r\in\R^J$ which dominate
the $J$ perspective cost sums of some feasible $(\alpha,y)$.
It is convex and upward closed; its intersection with every bounded
box is closed, by compactness of the feasible set and continuity of
the costs. If $t<\Gamma_{\rm pre}$, the vector $t\one$ is not in
$K$. A separating linear functional must have nonnegative coordinates,
since $K$ is upward closed. Normalize them to sum to one, obtaining
$\lambda\in\Delta_J$. Separation gives
\[
 t<\min_{\alpha,x^u}\sum_u\alpha_u
                         \sum_w\lambda_{j(w)}\phi_w(x^u_w)
       =\min_u m_u(\lambda).
\]
The reverse inequality holds for every $\lambda$, by the maximum
bound for a convex combination of the level costs. Let $t$ increase
to $\Gamma_{\rm pre}$ and use compactness of $\Delta_J$ to obtain
\[
             \Gamma_{\rm pre}=\max_{\lambda\in\Delta_J}
                                             \min_u m_u(\lambda).
\]
The finite-dimensional convex minimax used in
\eqref{eq:v29-inner-dual} is the same separating-hyperplane principle
applied to affine supporting functions: separate the strict sublevel
of the convex epigraph from a compact convex feasible set. Equivalently,
approximate the compact set of supporting slopes by finite nets,
apply finite convex minimax, and pass to the limit by uniform continuity.
All sets here are finite dimensional and compact; no infinite-path or
asymptotic interchange is involved.

Finally the maximizations defining $m_u$ are over independent compact
sets for the finitely many $u$. They can be chosen simultaneously,
so $\min_u\max_{q^u}$ equals $\max_{(q^u)}\min_u$ in this particular
separated situation. Substituting \eqref{eq:v29-inner-dual} proves
\eqref{eq:v29-dual}. The zero-weight convention is continuous, since
the bounded conjugate is multiplied by that weight. Thus the displayed
maximum is attained.
\end{proof}

When $c(t)=\min\{1,a\sqrt t\}$ with $a\ge1$, the needed functions
are particularly explicit:
\begin{equation}\label{eq:v29-cubic-dual}
 \phi(x)=\frac{x^3}{3a^2},\qquad
 \widehat\phi(r)=
 \begin{cases}
  \dfrac{2a}{3}r^{3/2},&0\le r\le a^{-2},\\[3pt]
  r-\dfrac1{3a^2},&a^{-2}\le r\le1.
 \end{cases}
\end{equation}
These identities follow by maximizing $rx-x^3/(3a^2)$ on $[0,1]$.
At a rational support point $z$, the tangent has slope $z^2/a^2$ and
intercept $-2z^3/(3a^2)$. Rational $a^2$ therefore permits exact
rational certificates, without floating-point convex optimization.

\subsection{The decoder descriptor and graph specialization}
Return to the positive graph experiment. For an internal visited set
$S$ and a positive allocation $\ell$ let
\[
 a_{S,\ell}=M C_{S,\ell}(a)/D_{\delta(S),\ell},
 \qquad
 c_S(t)=\min\{1,\min_{\ell\ {\rm positive}}
                                      a_{S,\ell}t^{k/2}\}.
\]
The empty cut has capacity one. The first-block event and all constants
are those of \upref{Lemma}{lem:v26-first-block}; its proof supplies
\eqref{eq:v29-local} for every fixed independent coding seed. The
scheduler chooses its first vertex before receiving any tape entry.

\begin{corollary}[A quantitative graph converse]
\label{cor:v29-graph}
For every fixed graph experiment, calibration and integer $M\ge1$,
\begin{equation}\label{eq:v29-graph}
 \mathcal R^{\rm max}_{G,M}(a)\ge\mathcal R^{\rm av}_{G,M}(a)
 \ge\beta(a)\Gamma_{\rm pre}(c)
 \ge\beta(a)\Gamma(c)
 \ge\frac{\beta(a)}{v-1}\int_0^\infty[1-H_G(t;M,a)]_+\,dt.
\end{equation}
The scheduler may be given its full revealed past in the lower bounds.
The original decoder sees only its label and the visited-set query
descriptor. Every retained scalar and graph profile upper bound remains
unchanged.
\end{corollary}
\begin{proof}
Apply Theorem~\ref{thm:v29-pre} to every admissible controller with
its independent seeds fixed, and then take the infimum over one common
controller. The capacities and all three variational quantities are
controller independent. Proposition~\ref{prop:v29-separator-dominance}
and the terminal feasibility argument justify the last, infinite-integral
form. The comparison of the two risks follows from their definitions.
\end{proof}

There is also a quantitative statement for a decoder receiving additional
finite history information. Suppose that, at a fixed visited set $S$,
its additional descriptor has at most $K_S$ possible values. After fixing
independent seeds it has at most $M K_S$ prediction centres, even when
the descriptor was selected measurably from the observed history.
The same unconditioned density bound thus proves the corollary with
\begin{equation}\label{eq:v29-descriptor}
 c^{K}_S(t)=\min\{1,\min_{\ell\ {\rm positive}}
                             K_S a_{S,\ell}t^{k/2}\}.
\end{equation}
For the full finite order prefix one may take $K_S=|S|!$. These are
explicitly different capacities, not a transfer of the original
constants to a stronger decoder. Alternatively, the state graph may
be lifted to prefix vertices and one may retain $M$ centres at each
lifted vertex. Both are legitimate lower bounds; neither grants any
continuous history for free to the decoder.

All dependence on graph size, query recovery, priors and acquisition
charts remains visible through $\beta$, the capacities and the finite
flow problems. The corollary makes no claim of graph-size-uniform
constants or of sharp statistical constants for every experiment.

\section{Quantitative refinement in the two-trial experiment}
\label{sec:v29-evaluated}
We now evaluate the refinement using the complete acquired law already
derived in \upref{Proposition}{prop:v27-density}. The detector, command
law, two-trial horizon, four equiprobable queries and decoder are
unchanged. Thus the comparison concerns complete lower bounds for the
same experiment, not merely a comparison of event masses.

\subsection{The first choice and the number of centres}
The graph has two vertices and one edge. Its subset lattice has two
internal vertices at a single level, and the first choice is made before
any command is revealed. Suppose that a tape event $A$ has probability
$\beta$, its acquired mean $z$ has subprobability density at most $H$
on that event, and the affine recovery norm from query space to $z$ is
$L$. For fixed independent seeds, the first vertex is fixed and the
decoder has at most $M$ centres. Hence, with
\[
                         a_M=2LMH/\beta,
\]
the node small-ball capacity is $c(t)=\min\{1,a_M\sqrt t\}$.
In both applications below $a_M>1$ for every $M\ge1$.

\begin{proposition}[Exact variational values for two first choices]
\label{prop:v29-two-choices}
For these capacities,
\begin{equation}\label{eq:v29-two-choices}
 \Gamma(c)=\frac1{12a_M^2},\qquad
 \Gamma_{\rm pre}(c)=\frac1{3a_M^2}.
\end{equation}
Consequently the seedwise converse gives
\begin{equation}\label{eq:v29-one-step-density}
           \mathcal R^{\rm av}_{G,M}
                         \ge\frac{\beta^3}{12L^2H^2}\,M^{-2}.
\end{equation}
The same bound holds for the finite-prefix decoder relaxation.
\end{proposition}
\begin{proof}
An unanchored occupation is $(x,1-x)$, $0\le x\le1$.
Equation~\eqref{eq:v29-power-cost} gives the objective
$[x^3+(1-x)^3]/(3a_M^2)$, whose minimum is $1/(12a_M^2)$ at $x=1/2$.
After fixing the first choice, the occupation is instead $(1,0)$ or
$(0,1)$. Each has cost $1/(3a_M^2)$. A seed mixture of the two has
that same single-level cost. This proves \eqref{eq:v29-two-choices}.
Multiplication by $\beta$ gives \eqref{eq:v29-one-step-density}.
The first order prefix consists of just the already identified first
vertex, so the stronger prefix decoder supplies no further information
at this only checkpoint.

For an explicit dual check, in the unanchored problem use support
points $1/2,1/2$ in \eqref{eq:v29-cubic-dual}. The path cost is
$1/(4a_M^2)$ and the two conjugate terms sum to $1/(6a_M^2)$.
Their difference is $1/(12a_M^2)$. For an anchored problem use support
point one at the forced vertex and zero at the other; the path cost is
$1/a_M^2$ and the conjugate sum is $2/(3a_M^2)$.
Thus feasible dual values match the stated primal values exactly.
\end{proof}

\subsection{A comparison retaining every constant}
The query vector is
\[
 q(z)=\left(\frac34,\frac12+\frac z8,
                    \frac12-\frac z8,\frac14\right),
 \qquad \|q(z)-q(z')\|_q^2=\frac{(z-z')^2}{128}.
\]
The affine recovery $z=4(q_2-q_3)$ has norm $L=8\sqrt2$, so $L^2=128$.
The two possible first accepted reports give atoms of masses $5/16$ and
$3/16$. The remaining, failure-report subprobability has density $f$
on $[10/21,14/27]$, total mass $1/2$, and exact maximum $26$.
These are the unconditional, evidence-weighted quantities from the
retained acquired-law proposition; the latent prior is not being used
as though it were the acquired-statistic law.

\begin{theorem}[An all-command failure certificate]
\label{thm:v29-evaluated}
For every integer $M\ge1$, in the actual positive two-trial experiment,
\begin{equation}\label{eq:v29-explicit}
           \mathcal R^{\rm av}_{G,M}\ge
                                  \frac{1}{8306688}\,M^{-2}.
\end{equation}
This holds in both the original decoder model and its finite-prefix
relaxation. Requiring in addition the failure report on the independently
prepared completion part of the exploration tape gives, with the same
query norm and decoder, the certificate
\begin{equation}\label{eq:v29-completion}
           \mathcal R^{\rm av}_{G,M}\ge
                                  \frac{1}{16613376}\,M^{-2}.
\end{equation}
The first-block certificate is therefore exactly twice the corresponding
complete-word certificate, with all density and recovery costs included.
\end{theorem}
\begin{proof}
Take $A$ to be the first failure report with no restriction on its
command. Its mass is $\beta=1/2$ and its acquired-mean density is bounded
by $H=26$. These inequalities hold before choosing or conditioning on
any order. The event depends only on the tape and is independent of
the coding seed. Substituting in
\eqref{eq:v29-one-step-density} gives
\[
 \frac{\beta^3}{12L^2H^2}
    =\frac{(1/2)^3}{12\cdot128\cdot26^2}
    =\frac1{8306688}.
\]
This argument discards the nonnegative error contributions of the two
atoms; it does not assert that they vanish for finite $M$.

For the comparison, the completion command has the same independent
uniform law on $[1/4,3/4]^2$. At every latent value $t$ its averaged
failure likelihood is
\[
 \E_{x,y}[1-xk_0(t)-yk_1(t)]
              =1-\tfrac12(k_0(t)+k_1(t))=\tfrac12.
\]
Thus the completion-failure restriction multiplies the entire first-block
subprobability measure by $1/2$, not just its total mass. It replaces
$(\beta,H)$ by $(\beta/2,H/2)$ and leaves $a_M=2LMH/\beta$ unchanged.
Formula~\eqref{eq:v29-one-step-density} is multiplied by $1/2$.
This restriction concerns a potential exploration-tape event used in
the converse. No unobserved completion report is supplied to the
predictor, and no conditioning on a future query outcome is part of
the implemented rule.
\end{proof}

The earlier evaluated subevent restricted the first command to
$[1/4,3/8]\times[5/8,3/4]$. Its actual mass and density bound were
\[
 \beta_{\Omega}=\frac{35}{1024},\qquad
 H_{\Omega}=\frac{177957}{20480},\qquad L^2=128.
\]
Keeping exactly these data, but retaining the pre-acquisition constraint,
multiplies the previous separator certificate by four:
\begin{equation}\label{eq:v29-restricted}
 \frac{\beta_{\Omega}^3}{12L^2H_{\Omega}^2}\,M^{-2}
    =\frac{1071875}{3113159280132096}\,M^{-2}.
\end{equation}
This isolates the effect of fixing independent seeds before the initial
choice. Passing additionally to the complete first-failure density
yields \eqref{eq:v29-explicit}. Relative to the earlier, fully evaluated
restricted-event separator coefficient, the exact ratio is
\begin{equation}\label{eq:v29-improvement}
 \frac{1/8306688}{1071875/12452637120528384}
                    =\frac{1499109768}{1071875}>1398.
\end{equation}
This larger ratio combines two distinct changes: the factor four from
pre-acquisition information, and use of the entire known first-failure
law. It is not attributed to a larger event mass alone.

\subsection{Relation to the sharp statistical coefficient}
The exact reduction and high-resolution theorem retained from the
preceding section give
\begin{equation}\label{eq:v29-sharp-retained}
 \mathcal R^{\rm av}_{G,M}=\frac{D_M(\nu)}{128},\qquad
 \lim_{M\to\infty}M^2\mathcal R^{\rm av}_{G,M}
  =\kappa_{\rm ex}
  =\frac1{1536}\left(\int_{10/21}^{14/27}f(z)^{1/3}\,dz\right)^3.
\end{equation}
In particular, the asymptotic ratio of the statistical optimum to the
new lower certificate is exactly
\begin{equation}\label{eq:v29-gap}
                  8306688\,\kappa_{\rm ex}.
\end{equation}
The rational endpoint-sum enclosure for the retained integral gives
$4.8846\,10^{-7}<\kappa_{\rm ex}<4.8871\,10^{-7}$, and hence
\[
                     4.05<8306688\,\kappa_{\rm ex}<4.06.
\]
The exact integral, rather than its displayed decimal enclosure, defines
the coefficient. The bound \eqref{eq:v29-explicit} is valid for every
integer budget. The comparison by $4.06$ concerns the high-resolution
limit and is not a finite-budget relative-error assertion. Nor is an
equality with the maximum-tape criterion asserted.

The example separates three quantities: the separator certificate, the
stronger pre-acquisition occupation certificate, and the sharp statistical
risk. The first two arise from specified local information; the last
uses the full mixed acquired law. The general flow duality characterizes
the first kind of variational information exactly, while the statistical
constant in this example is fixed by the retained scalar quantization
argument. These statements complement one another and do not identify
the two optimization problems.


\appendix
\section{The acquired law in the evaluated experiment}
\label{app:v29-density}
For completeness we record the calculation from the retained v27 evaluated-model module. The prior is uniform on $[0,1]$, the cells are
\[
 k_0(t)=\frac12+\frac t4,\qquad k_1(t)=\frac12-\frac t4,
\]
and the acquisition command $(x,y)$ has density $4$ on $[1/4,3/4]^2$. The first-report likelihoods are $xk_0$, $yk_1$ and $1-xk_0-yk_1$. After an accepted report the command scalar cancels from Bayes' formula. The corresponding posterior means and unconditional masses are
\[
 (z_0,p_0)=(8/15,5/16),\qquad (z_1,p_1)=(4/9,3/16).
\]
On failure put
\[
 E(x,y)=1-\frac58x-\frac38y,
 \qquad z=\frac12+\frac{y-x}{48E(x,y)},
 \qquad u=48(z-1/2).
\]
The failure-weighted input density is $4E(x,y)$, not the unweighted command density. The inverse change of variables is
\[
 y=\frac{8u+(8-5u)x}{8+3u},\qquad
 E=\frac{8(1-x)}{8+3u},\qquad
 \frac{\partial y}{\partial z}=\frac{3072(1-x)}{(8+3u)^2}.
\]
Thus the failure-weighted joint density in $(x,z)$ is
$98304(1-x)^2/(8+3u)^3$. At fixed $u$ the admissible $x$ intervals are
\[
 \begin{cases}
 [(8-29u)/(4(8-5u)),\ 3/4],&-8/7\le u\le0,\\
 [1/4,\ (24-23u)/(4(8-5u))],&0\le u\le8/9.
 \end{cases}
\]
Integrating in $x$ gives the continuous subprobability density
\begin{equation}\label{eq:v29-density-explicit}
 f(z)=\begin{cases}
 512\left(\dfrac{27}{(8-5u)^3}-\dfrac1{(8+3u)^3}\right),&-8/7\le u\le0,\\[5pt]
 512\left(\dfrac{27}{(8+3u)^3}-\dfrac1{(8-5u)^3}\right),&0\le u\le8/9,\\[4pt]
 0,&\text{otherwise}.
 \end{cases}
\end{equation}
Its support is $[10/21,14/27]$, both endpoint values are zero, and its value at $u=0$ is $26$. On the negative interval its $u$ derivative is
\[
 512\left(\frac{405}{(8-5u)^4}+\frac9{(8+3u)^4}\right)>0,
\]
whereas on the positive interval it is
\[
 -512\left(\frac{243}{(8+3u)^4}+\frac{15}{(8-5u)^4}\right)<0.
\]
The total failure mass is the integral of $4E$ on the command square, namely $1/2$. Consequently the full law is
\[
 \nu=\frac5{16}\delta_{8/15}+\frac3{16}\delta_{4/9}+f(z)\,dz.
\]
The four equiprobable corner queries have the vector $q(z)$ displayed in Section~\ref{sec:v29-evaluated}. Orthogonal projection of arbitrary prediction centres onto that affine line, and nearest-centre assignment, prove the exact reduction $\mathcal R^{\mathrm{av}}_{G,M}=D_M(\nu)/128$. Conversely, one transition computes $z$, stores the nearest-centre index and uses $q$ at decoding. The first vertex was chosen independently of the tape, so independent randomization cannot lower this scalar infimum.

The sharp continuous-density quantization coefficient is classical \cite{GrafLuschgy2000}; its complete interval-partition proof, including finite atoms, is retained in v27. Briefly, for a fixed partition into intervals of lengths $l_i$, let $h_i$ and $H_i$ be the infimum and supremum of $f$. Intersecting $M$ Voronoi cells with $L$ partition intervals produces at most $M+L-1$ pieces. The error on an interval of length $d$ is at least $h_i d^3/12$. Convexity and H\"older's inequality give the lower bound
\[
 D_M(f\,dz)\ge
       \frac{(\sum_i h_i^{1/3}l_i)^3}{12(M+L-1)^2}.
\]
For the upper bound, assign $m_i$ midpoint cells to interval $i$, with $m_i/M$ converging to $H_i^{1/3}l_i/\sum_r H_r^{1/3}l_r$. Their assigned error is at most $\sum_i H_i l_i^3/(12m_i^2)$. Let $M$ grow at fixed partition, then refine the partition; continuity of $f^{1/3}$ makes the bounds coincide. Zero-density intervals contribute nothing. Adding a centre for each of the two atoms uses only two labels and changes no leading coefficient. This yields \eqref{eq:v29-sharp-retained}.

The rational enclosure in the accompanying diagnostic does not use numerical quadrature as a proof. On each of the two monotone intervals, take $8192$ subintervals in $u$, and use $dz=du/48$. At every rational endpoint the value of $f$ is rational. Bisection of integers bounds its cube root between consecutive multiples of $10^{-12}$. Lower and upper endpoint sums then enclose the integral. Cubing and dividing by $1536$ gives the bounds used in Section~\ref{sec:v29-evaluated}.

\begin{thebibliography}{9}
\bibitem{Rockafellar1964}
R.~T.~Rockafellar,
Minimax theorems and conjugate saddle-functions,
\emph{Mathematica Scandinavica} \textbf{14} (1964), 151--173.
\bibitem{GrafLuschgy2000}
S.~Graf and H.~Luschgy,
\emph{Foundations of Quantization for Probability Distributions},
Lecture Notes in Mathematics \textbf{1730}, Springer-Verlag, Berlin, 2000.
\end{thebibliography}
\end{document}
